\section{Introduction}\label{sec:intro}

Metaheuristics remain the workhorse for large routing problems, and a growing literature embeds machine-learning components in them to select search operators, tune parameters or guide moves \citep{karimimamaghan2022,bengio2021}. A common design is \emph{adaptive operator selection} (AOS): a policy observes the success of the operators applied so far and shifts the probability of choosing each one. Multi-armed-bandit policies are a popular choice for this role \citep{dacosta2008,fialho2010,lagos2024}, and adaptive weights are a standard ingredient of adaptive large neighborhood search \citep{ropke2006}. Newer proposals make operator selection a learned decision \citep{karimimamaghan2023,pei2023,xu2025}.

Whether such learning \emph{helps} is, however, a question that is rarely isolated. A learned policy is always evaluated inside a search with its own stopping rule, acceptance rule and operator order, and gains attributed to the policy may belong to the structure around it. The one meta-analytic estimate we know of is modest: across 25 implementations of adaptive large neighborhood search, the adaptive layer improved the objective by 0.14\% on average (95\% interval 0.06--0.21\%) \citep{turkes2021}.

This paper audits the learning component of one recent, representative method. \citet{kalatzantonakis2023} couple a variable neighborhood search (VNS) \citep{mladenovic1997,hansen2001} for the capacitated vehicle routing problem (CVRP) with a hybrid upper-confidence-bound and gradient-bandit policy (HG-UCB) and report that it improves on plain VNS by 0.29 percentage points (pp) of optimality gap on average over three benchmark sets. We ask five research questions.
\begin{description}
\item[RQ1 (reproduction)] Does the published improvement replicate, and what in the reward design determines the policy's behavior?
\item[RQ2 (structure)] How large are the effects of structural choices of the search (inner-loop termination, acceptance, operator order, operator pool) relative to learning?
\item[RQ3 (deterministic operators)] Under an equal budget, does adaptive selection beat non-adaptive (uniform) selection, and does the answer survive tuning, further policy families, larger instances, an efficient implementation of the costliest operator, a state-of-the-art solver and a positive control that shows the protocol can see a gain?
\item[RQ4 (stochastic operators)] With the stochastic operators of an adaptive large neighborhood search (CVRP and CVRP with time windows), does adaptive selection help, and is the gain adaptation or merely a good fixed weighting of the operators?
\item[RQ5 (generality)] Do the findings hold for another problem (the traveling salesman problem) and for one with new constraints (the CVRP with time windows), and where do they not?
\end{description}

We answer them with sixteen studies, each pre-specified in a dated research log before its data were produced, using 74 CVRP benchmark instances, 32 TSPLIB instances, 56 Solomon instances and 164 synthetic ones, ten random seeds (five on most synthetic sets), budgets measured in operator applications, in distance evaluations or in generations, paired nonparametric tests with Holm correction, bootstrap intervals, equivalence tests, a positive control and a state-of-the-art reference solver.

\paragraph{Contributions}
\begin{enumerate}
\item A reproduction of HG-UCB-guided VNS with its interpretation gaps documented. The published improvement is not reproduced, and we identify a failure mechanism: a reward in raw cost units saturates the policy's softmax, and normalizing it removes a loss of \num{s1a_hgucb_literal_hgucb_norm_800} to \num{s1a_hgucb_literal_hgucb_norm_4000} pp (Sections~\ref{sec:repro} and~\ref{sec:r1}).
\item Evidence, under matched budgets and with equivalence tests, that in the CVRP none of ten adaptive policy families beats uniform random selection with deterministic local-search operators (\num{u_n} contrasts, none better), and that the search structure (termination, operator order, operator pool, ruin and recreate) moves the gap far more (Sections~\ref{sec:r2} and~\ref{sec:r3}).
\item A positive control, repeated in three problems, showing that the protocol detects adaptive gains when they exist (Sections~\ref{sec:r4} and~\ref{sec:r6}).
\item A regularity across three routing problems: adaptive selection beats uniform selection where the pool contains costly operators that are not worth their cost (the larger pool is worse than the smaller one: the TSP, time windows, planted operators), and a non-adaptive design (a smaller pool, a descent ordered by cost, a tuned mixture or a fixed repair) obtains the same gain; in no setting is a policy significantly better than the best of these alternatives (the one policy that beats any fixed design, HG-UCB with time windows, beats a single repair that is no better than uniform selection, not the tuned mixture). A prospective test on eight further nested pools confirms the second statement without exception (\num{s15_foc_nb} of \num{s15_foc_n} contrasts against the fixed-order descent of the same pool favor a policy) but supports the first only as a tendency (Sections~\ref{sec:r5}--\ref{sec:r7}).
\item Practical guidance for choosing the selection layer of a local search, and an open, auditable protocol: pre-specification entries, exact-regression controls, operators validated by brute force, tables generated from raw results and recomputed independently (Section~\ref{sec:discussion}).
\end{enumerate}
